\section{Một phép đối chiếu đã được làm}
\label{sec:ch12-mot-phep-doi-chieu}

\index{đối chiếu chỉ số với người|(}%
Cuộc kiểm đếm cho biết mũi tên nằm ngang của
hình~\ref{fig:ch12-hai-nhanh-danh-gia} hầu như không được vẽ, nhưng nó không
cho biết chuyện gì xảy ra ở những lần hiếm hoi mà nó được vẽ. Liedeker, Ell,
Cimiano và Düsing, bốn tác giả ở Đại học Bielefeld, chạy đúng phép đối chiếu
ấy trong một nghiên cứu công bố tháng 3 năm 2026 và được nhận ở
hội nghị World Conference on eXplainable Artificial Intelligence 2026, cùng
loạt hội nghị đã nhận khảo sát mà chương~\ref{ch:lime-nao-dang-tin}
đọc~\autocite{p25userperception}. Đây là bài đầu tiên trong danh mục của sách
đặt điểm số của một chỉ số tự động cạnh đánh giá của người trên cùng những lời
giải thích rồi báo cáo mức trùng khớp.

\index{lời giải thích phản thực}%
Đối tượng của nó phải nói ngay, vì nó không phải đối tượng của Phần~II. Bài
báo chấm các chỉ số của \tn{lời giải thích phản thực}{counterfactual
explanation}, tức một họ phương pháp giải thích khác hẳn attribution đặc trưng:
thay vì gán một điểm số cho từng \tn{đặc trưng}{feature}, một lời giải thích phản thực trả
về một phiên bản đã sửa của chính đầu vào, sửa ít nhất có thể, sao cho mô hình
đổi dự đoán sang lớp khác. Người đọc nhìn vào đó thì biết cần đổi gì để kết
quả khác đi.

Chữ \enquote{phản thực} ở đây không cùng nghĩa với chữ \tn{phản thực}{counterfactual} mà Phần~IV đã
dùng, và hai nghĩa ấy phải giữ tách nhau. Từ chương~\ref{ch:khung-hop-nhat}
trở đi, phép thử phản thực là thao tác bỏ đi hoặc thay đổi những phần tử được
chấm điểm cao rồi xem đầu ra của mô hình đổi bao nhiêu, tức một cách đo một
lời giải thích đã có. Ở chương này, một lời giải thích phản thực là bản thân
lời giải thích. Cùng một chữ, một bên là dụng cụ đo và một bên là vật được đo.
Đây là trường hợp cùng loại với chỗ va chữ mà
chương~\ref{ch:lime-nao-dang-tin} đã phải gỡ, chỉ khác là ở đó hai bài báo va
nhau ở chữ fidelity còn ở đây chúng va nhau ở chữ phản thực.

Hệ quả của ranh giới ấy giống hệt hệ quả mà
mục~\ref{sec:ch09-doi-tuong-khac-nhau} rút ra cho bài báo neo. Bài báo neo đo
các chỉ số của \tn{chuỗi suy luận}{chain-of-thought}, bài này đo các chỉ số của lời giải thích phản
thực, và không bài nào đo các họ chỉ số của
chương~\ref{ch:do-do-trung-thuc}. Cái đi qua được ranh giới lần này là gì thì
mục~\ref{sec:ch12-ket-qua-va-ranh-gioi} nói, sau khi đã có kết quả để nói về.

\index{đối chiếu chỉ số với người!thiết kế nghiên cứu}%
Thiết kế nghiên cứu gồm ba phần, và phần đầu là sinh lời giải thích. Nó chạy
trên ba bộ dữ liệu bảng lấy từ kho UCI, chọn theo tiêu chí đặc
trưng và nhãn phải trực giác đủ để người không chuyên đọc được: Mushroom, phân
loại nấm ăn được hay độc, Obesity Levels, phân loại mức béo phì thành bảy lớp,
và Heart Disease, dự đoán có bệnh tim hay không. Mô hình nền là XGBoost, đạt
điểm F1 từ 0.85 trở lên trên cả ba. Lời giải thích phản thực được sinh cho mọi
điểm trong tập kiểm thử bằng phương pháp Counterfactuals Guided by Prototypes,
qua thư viện Alibi Explain; một lời giải thích được coi là hợp lệ khi dự đoán
của mô hình trên điểm mới sinh đã lật sang lớp
đích~\autocite{p25userperception}.

Từ khối lời giải thích hợp lệ ấy, các tác giả rút một mẫu nhỏ để đưa vào
nghiên cứu người dùng, và cách rút mẫu đáng ghi vì nó quyết định phổ của cái
được chấm. Họ phân cụm bằng k-means trên chính vector bảy chỉ số của từng lời
giải thích, chia số suất theo cỡ cụm sao cho mỗi cụm góp ít nhất một lời giải
thích, rồi rút ngẫu nhiên đều trong từng cụm. Kết quả là 85 lời giải thích,
chia 30 cho Mushroom, 30 cho Obesity Levels và 25 cho Heart
Disease~\autocite{p25userperception}. Cách rút này giữ cho mẫu phủ được cấu
trúc của toàn bộ không gian chỉ số thay vì dồn vào vùng đông nhất; hệ quả là
tương quan tính về sau không bị thu hẹp chỉ vì các lời giải thích được chấm
quá giống nhau, và hệ quả ấy là điều sách đọc ra chứ không phải điều bài báo
nêu.

\index{đối chiếu chỉ số với người!nghiên cứu người dùng}%
Phần thứ hai là thu đánh giá của người, và nó chạy trên nền tảng Prolific. Người tham gia được giới thiệu
về lời giải thích phản thực bằng một ví dụ trên bộ dữ liệu tín dụng Đức đã
được giản lược, rồi mỗi người nhận một trong 20 gói, mỗi gói 12 lời giải
thích, bốn cho mỗi bộ dữ liệu. Thứ tự các bộ dữ liệu được xáo, còn các lời
giải thích cùng một bộ thì đặt liền nhau để giảm tải nhận thức. Mỗi lời giải
thích hiện ra dưới dạng bảng đối chiếu điểm gốc với phiên bản đã sửa, các đặc
trưng bị đổi được tô sáng, kèm một đoạn mô tả ngắn bằng
lời~\autocite{p25userperception}.

Mỗi lời giải thích được chấm theo năm chiều, mỗi chiều là một câu hỏi trên
\tn{thang Likert}{Likert scale} bốn mức, từ 1 là \enquote{Definitely Yes} tới
4 là \enquote{Definitely No}, cộng thêm lựa chọn không biết. Năm chiều là: lớp
mà mô hình dự đoán có chính xác không; lời giải thích có hiểu được không; lời
giải thích có hợp lý không; lời giải thích có đủ chi tiết không; và lời giải
thích có làm người đọc hài lòng không~\autocite{p25userperception}. Cần chú ý
chiều của thang: 1 là câu trả lời tích cực nhất, nên số càng nhỏ càng tốt.

Đọc năm chiều ấy cạnh cặp khái niệm của
mục~\ref{sec:ch01-do-trung-thuc-va-tinh-hop-ly} thì thấy cả năm đều nằm ở phía
tính hợp lý. Không chiều nào hỏi người tham gia rằng lời giải thích có phản
ánh đúng mô hình hay không, và cũng không thể hỏi, vì người tham gia không
nhìn thấy mô hình. Nhận xét này là của sách, và nó không phải một lời chê
thiết kế: nó là điều mà mũi tên chấm vòng lên phía mô hình trong
hình~\ref{fig:ch12-hai-nhanh-danh-gia} đã vẽ sẵn.

\index{đối chiếu chỉ số với người!cỡ mẫu}%
Mẫu người cuối cùng gồm 167 người, sau khi một người trượt các câu kiểm tra
độ chú ý và bị loại. Tuổi từ 20 tới 72, trung bình 40.85 với độ lệch chuẩn
13.05; 77.8\% có bằng đại học trở lên; kinh nghiệm học máy tự khai trung bình
2.67 với độ lệch chuẩn 0.93 trên thang bốn mức. Tổng cộng 2\,004 lượt chấm được
thu, tức trung bình 23.58 bộ chấm đầy đủ cho mỗi lời giải
thích~\autocite{p25userperception}.

Năm chiều ấy về sau được gộp thành một điểm duy nhất, và bước gộp cần được
biện minh chứ không được làm mặc nhiên, vì mọi kết quả của
mục~\ref{sec:ch12-ket-qua-va-ranh-gioi} tính trên điểm gộp ấy. Các tác giả
biện minh bằng hai phép kiểm, mỗi phép trả lời một câu hỏi khác nhau.

Câu hỏi thứ nhất là những người chấm có đồng thuận với nhau không. Nếu hai
người cùng đọc một lời giải thích mà cho hai điểm rất khác nhau thì điểm ấy
nói về người chấm nhiều hơn nói về lời giải thích. Các tác giả đo bằng hệ số
tương quan nội lớp, tức tỉ lệ phương sai của các điểm chấm sinh ra từ khác
biệt giữa các lời giải thích chứ không từ khác biệt giữa những người chấm. Kết
quả là thấp khi xét từng người chấm một, và cao hơn nhưng vẫn chỉ ở mức vừa
phải khi lấy trung bình nhiều người, đúng như các nghiên cứu người dùng trước
đó trong XAI~\autocite{p25userperception}. Bài báo không in giá trị cụ thể
nào, nên sách không in giá trị nào.

Câu hỏi thứ hai là năm chiều có đang đo cùng một thứ không, và câu này mới là
câu cho phép gộp. Hai phép kiểm trả lời nó. Hệ số Cronbach alpha đo mức năm
chiều dao động cùng nhau qua các lời giải thích, và giá trị 0.88 nằm ở mức
được coi là cao. Phân tích thành phần chính đi xa hơn một bước: nó tìm trục mà
theo đó năm chiều biến thiên nhiều nhất, và ở đây trục thứ nhất một mình đã
giải thích 74.1\% phương sai của cả năm, nghĩa là gần ba phần tư những gì năm
chiều phân biệt được đều nằm trên một trục duy
nhất~\autocite{p25userperception}. Với hai kết quả đó, lấy trung bình năm
chiều không làm mất mấy thông tin, nên các tác giả gộp chúng thành một điểm
duy nhất mà họ đặt tên là Combined Quality Score, và mọi phân tích về sau chạy
trên điểm gộp ấy bên cạnh năm chiều riêng. Trên toàn bộ 85 lời giải thích,
điểm gộp trung bình là 2.10 với độ lệch chuẩn 0.22, tức nằm giữa mức 2 và mức 3
của thang bốn mức. Bài báo không đọc gì thêm từ con số ấy: nó dùng điểm gộp làm
biến phụ thuộc cho các phân tích sau chứ không bình luận mức chất lượng mà 2.10
biểu thị, nên sách cũng không bình luận thay.

\index{đối chiếu chỉ số với người!lực kiểm định}%
Còn một phần nữa của thiết kế mà bài báo làm và phần lớn nghiên cứu người dùng
không làm, và chính phần này quyết định kết quả ở
mục~\ref{sec:ch12-ket-qua-va-ranh-gioi} được đọc như thế nào. Đó là phân tích
\tn{lực kiểm định}{statistical power}: tính trước xem với cỡ mẫu đang có, hiệu
ứng nhỏ tới đâu thì vẫn phát hiện được, và nhỏ hơn ngưỡng đó thì không. Lực
kiểm định là xác suất phép thử báo có hiệu ứng khi hiệu ứng thật sự tồn tại ở
một cỡ cho trước, nên nói một nghiên cứu có lực 80\% với hiệu ứng cỡ ấy là nói
rằng nếu hiệu ứng đúng cỡ ấy có thật thì cứ mười lần chạy sẽ có tám lần thấy
được nó. Với mức ý nghĩa 0.05 và cỡ mẫu 25 tới 30 lời giải thích cho mỗi bộ dữ
liệu, nghiên cứu đạt lực 80\% cho các hiệu ứng lớn, tức hệ số tương quan từ
0.50 trở lên. Với phần hồi quy dùng bảy chỉ số làm biến dự báo, cỡ ấy đủ để phát
hiện hệ số xác định $R^2$ từ 0.40 trở lên; muốn phát hiện $R^2$ bằng 0.30 thì
cần khoảng 42 lời giải thích cho mỗi bộ dữ liệu, và $R^2$ bằng 0.15 thì cần
khoảng 89~\autocite{p25userperception}.

Các tác giả nói rõ mình chọn cỡ ấy vì lý do gì, và lý do đó là một lập luận về
mục đích chứ không phải một lời biện hộ. Câu hỏi trung tâm của nghiên cứu là
liệu các chỉ số tự động có làm được đại lượng thay thế cho đánh giá của người
hay không. Nếu một chỉ số, hoặc một tổ hợp chỉ số, xấp xỉ được đánh giá của
người thì hiệu ứng phải lớn, và cỡ mẫu này thừa sức thấy. Còn một chỉ số chỉ
tương quan yếu với cảm nhận của người, dưới 0.30, thì các tác giả cho rằng nó
ít có giá trị thực dụng~\autocite{p25userperception}. Nghiên cứu vì thế được
dựng để trả lời câu hỏi \enquote{có dùng thay được không}, và nó thừa nhận
rằng câu hỏi \enquote{có tương quan chút nào không} thì cần mẫu lớn hơn.
